\documentclass[conference]{IEEEtran}
\usepackage[T1]{fontenc}
\usepackage{amsmath,amssymb,amsthm}
\usepackage{booktabs}
\usepackage{listings}
\usepackage{xcolor}
\usepackage{url}

\newtheorem{theorem}{Theorem}
\newtheorem{lemma}{Lemma}
\newtheorem{corollary}{Corollary}

% Every measured number comes from eval/results/ through eval/paper_numbers.py.
% Generated by eval/paper_numbers.py from eval/results/. Do not edit.
\newcommand{\RqOneMsMin}{22}
\newcommand{\RqOneMsMax}{39}
\newcommand{\RqOneMsMed}{29}
\newcommand{\TabOpt}{1.49}
\newcommand{\TabOz}{1.81}
\newcommand{\TabBase}{1.82}
\newcommand{\TabSpeedup}{1.21}
\newcommand{\TabStart}{1.02}
\newcommand{\TabProcOpt}{0.47}
\newcommand{\TabProcOz}{0.79}
\newcommand{\TabCut}{40}
\newcommand{\ImgPl}{3.52}
\newcommand{\ImgPlMem}{10}
\newcommand{\ImgBase}{4.01}
\newcommand{\ImgBaseMem}{922}
\newcommand{\ImgHw}{3.54}
\newcommand{\ImgHwMem}{12}
\newcommand{\FusExact}{5.56}
\newcommand{\FusApprox}{5.26}
\newcommand{\FusSave}{5}
\newcommand{\FusPsnr}{57.0}
\newcommand{\FusPsnrMin}{51.8}
\newcommand{\CsvTenK}{0.01}
\newcommand{\CsvTenKRead}{0.06}
\newcommand{\CsvHundredK}{0.03}
\newcommand{\CsvHundredKRead}{0.08}
\newcommand{\CsvMillion}{0.13}
\newcommand{\CsvMillionRead}{0.62}
\newcommand{\ImgHundred}{0.02}
\newcommand{\ImgHundredRead}{0.14}
\newcommand{\ImgThousand}{0.08}
\newcommand{\ImgThousandRead}{1.48}
\newcommand{\ImgThreeK}{0.33}
\newcommand{\ImgThreeKRead}{5.73}
\newcommand{\CsvRatio}{0.2}
\newcommand{\CsvSpeedup}{4.7}
\newcommand{\CsvArrowRead}{0.11}
\newcommand{\CsvArrow}{0.8}
\newcommand{\HdrRatio}{17}
\newcommand{\PqOne}{0.015}
\newcommand{\PqOneRead}{0.66}
\newcommand{\PqOneSize}{22}
\newcommand{\PqFour}{0.016}
\newcommand{\PqFourRead}{0.19}
\newcommand{\PqFourSize}{87}
\newcommand{\PqSixteen}{0.018}
\newcommand{\PqSixteenRead}{0.84}
\newcommand{\PqSixteenSize}{346}
\newcommand{\ScTabOkPl}{0.02}
\newcommand{\ScTabOkScript}{4.51}
\newcommand{\ScTabOkRun}{4.74}
\newcommand{\ScTabNamePl}{0.02}
\newcommand{\ScTabNameScript}{2.81}
\newcommand{\ScTabContractPl}{0.02}
\newcommand{\ScTabContractScript}{1.79}
\newcommand{\ScImgOkPl}{0.33}
\newcommand{\ScImgOkScript}{8.53}
\newcommand{\ScImgOkRun}{10.51}
\newcommand{\ScImgBrokenPl}{0.26}
\newcommand{\ScImgBrokenScript}{9.15}
\newcommand{\ScImgFloatPl}{0.26}
\newcommand{\ScImgFloatScript}{8.54}
\newcommand{\RbTotal}{44}
\newcommand{\RbEarly}{18}
\newcommand{\RbEarlyPct}{41}
\newcommand{\RbCompile}{12}
\newcommand{\RbContract}{6}
\newcommand{\RbDesign}{3}
\newcommand{\RbNone}{23}
\newcommand{\RbBefore}{15}
\newcommand{\RbBeforePct}{34}
\newcommand{\RbCases}{21}
\newcommand{\RbPass}{21}
\newcommand{\RbFormat}{6}
\newcommand{\RbOperation}{5}
\newcommand{\RbDates}{4}
\newcommand{\RbLabels}{3}
\newcommand{\RbProduced}{2}
\newcommand{\RbContractMissing}{2}
\newcommand{\RbModel}{1}
\newcommand{\CertBatches}{244}
\newcommand{\CertDeopted}{35}
\newcommand{\CertRescued}{35}
\newcommand{\CertFpCorrect}{115}
\newcommand{\CertFpNeedless}{129}
\newcommand{\CertFpUnsound}{0}
\newcommand{\CertFpCorrectPct}{47}
\newcommand{\CertFpNeedlessPct}{53}
\newcommand{\CertFactsCorrect}{122}
\newcommand{\CertFactsNeedless}{81}
\newcommand{\CertFactsUnsound}{41}
\newcommand{\CertFactsCorrectPct}{50}
\newcommand{\CertFactsNeedlessPct}{33}
\newcommand{\CertPanderaAllCorrect}{121}
\newcommand{\CertPanderaAllNeedless}{82}
\newcommand{\CertPanderaAllUnsound}{41}
\newcommand{\CertPanderaAllCorrectPct}{50}
\newcommand{\CertPanderaAllNeedlessPct}{34}
\newcommand{\CertPanderaReadCorrect}{124}
\newcommand{\CertPanderaReadNeedless}{79}
\newcommand{\CertPanderaReadUnsound}{41}
\newcommand{\CertPanderaReadCorrectPct}{51}
\newcommand{\CertPanderaReadNeedlessPct}{32}
\newcommand{\CertGxAllCorrect}{138}
\newcommand{\CertGxAllNeedless}{106}
\newcommand{\CertGxAllUnsound}{0}
\newcommand{\CertGxAllCorrectPct}{57}
\newcommand{\CertGxAllNeedlessPct}{43}
\newcommand{\CertGxReadCorrect}{140}
\newcommand{\CertGxReadNeedless}{104}
\newcommand{\CertGxReadUnsound}{0}
\newcommand{\CertGxReadCorrectPct}{57}
\newcommand{\CertGxReadNeedlessPct}{43}
\newcommand{\CertTfdvAllCorrect}{147}
\newcommand{\CertTfdvAllNeedless}{31}
\newcommand{\CertTfdvAllUnsound}{66}
\newcommand{\CertTfdvAllCorrectPct}{60}
\newcommand{\CertTfdvAllNeedlessPct}{13}
\newcommand{\CertTfdvReadCorrect}{150}
\newcommand{\CertTfdvReadNeedless}{28}
\newcommand{\CertTfdvReadUnsound}{66}
\newcommand{\CertTfdvReadCorrectPct}{61}
\newcommand{\CertTfdvReadNeedlessPct}{11}
\newcommand{\CertFullCorrect}{244}
\newcommand{\CertFullNeedless}{0}
\newcommand{\CertFullUnsound}{0}
\newcommand{\CertFullCorrectPct}{100}
\newcommand{\CertFullNeedlessPct}{0}
\newcommand{\CertDeoptCorrect}{244}
\newcommand{\CertDeoptNeedless}{0}
\newcommand{\CertDeoptUnsound}{0}
\newcommand{\CertDeoptCorrectPct}{100}
\newcommand{\CertDeoptNeedlessPct}{0}
\newcommand{\RealRuns}{238}
\newcommand{\RealDeopted}{9}
\newcommand{\RealRescued}{9}
\newcommand{\RealFreshFails}{17}
\newcommand{\RealTypeRefusals}{24}
\newcommand{\RealFpCorrect}{17}
\newcommand{\RealFpNeedless}{221}
\newcommand{\RealFpUnsound}{0}
\newcommand{\RealFpCorrectPct}{7}
\newcommand{\RealFpNeedlessPct}{93}
\newcommand{\RealFactsCorrect}{19}
\newcommand{\RealFactsNeedless}{219}
\newcommand{\RealFactsUnsound}{0}
\newcommand{\RealFactsCorrectPct}{8}
\newcommand{\RealFactsNeedlessPct}{92}
\newcommand{\RealPanderaAllCorrect}{19}
\newcommand{\RealPanderaAllNeedless}{219}
\newcommand{\RealPanderaAllUnsound}{0}
\newcommand{\RealPanderaAllCorrectPct}{8}
\newcommand{\RealPanderaAllNeedlessPct}{92}
\newcommand{\RealPanderaReadCorrect}{19}
\newcommand{\RealPanderaReadNeedless}{219}
\newcommand{\RealPanderaReadUnsound}{0}
\newcommand{\RealPanderaReadCorrectPct}{8}
\newcommand{\RealPanderaReadNeedlessPct}{92}
\newcommand{\RealGxAllCorrect}{17}
\newcommand{\RealGxAllNeedless}{221}
\newcommand{\RealGxAllUnsound}{0}
\newcommand{\RealGxAllCorrectPct}{7}
\newcommand{\RealGxAllNeedlessPct}{93}
\newcommand{\RealGxReadCorrect}{17}
\newcommand{\RealGxReadNeedless}{221}
\newcommand{\RealGxReadUnsound}{0}
\newcommand{\RealGxReadCorrectPct}{7}
\newcommand{\RealGxReadNeedlessPct}{93}
\newcommand{\RealTfdvAllCorrect}{51}
\newcommand{\RealTfdvAllNeedless}{163}
\newcommand{\RealTfdvAllUnsound}{24}
\newcommand{\RealTfdvAllCorrectPct}{21}
\newcommand{\RealTfdvAllNeedlessPct}{68}
\newcommand{\RealTfdvReadCorrect}{138}
\newcommand{\RealTfdvReadNeedless}{73}
\newcommand{\RealTfdvReadUnsound}{27}
\newcommand{\RealTfdvReadCorrectPct}{58}
\newcommand{\RealTfdvReadNeedlessPct}{31}
\newcommand{\RealFullCorrect}{215}
\newcommand{\RealFullNeedless}{23}
\newcommand{\RealFullUnsound}{0}
\newcommand{\RealFullCorrectPct}{90}
\newcommand{\RealFullNeedlessPct}{10}
\newcommand{\RealDeoptCorrect}{221}
\newcommand{\RealDeoptNeedless}{17}
\newcommand{\RealDeoptUnsound}{0}
\newcommand{\RealDeoptCorrectPct}{93}
\newcommand{\RealDeoptNeedlessPct}{7}
\newcommand{\RealValidatorStopMin}{92}
\newcommand{\RealValidatorStopMax}{93}
\newcommand{\RealValidatorWrongMax}{27}
\newcommand{\CorpusAll}{372}
\newcommand{\CorpusTest}{272}
\newcommand{\CorpusLifted}{130}
\newcommand{\CorpusCompiled}{124}
\newcommand{\CorpusStatements}{189}
\newcommand{\CorpusChanges}{328}
\newcommand{\CorpusComplete}{30}
\newcommand{\CorpusRuns}{1364}
\newcommand{\CorpusAdmit}{1357}
\newcommand{\CorpusAdmitPct}{99}
\newcommand{\CorpusSame}{1357}
\newcommand{\CorpusUnsound}{0}
\newcommand{\CorpusStopped}{0}
\newcommand{\CorpusStoppedRuns}{0}
\newcommand{\CorpusRefusedCrash}{7}
\newcommand{\CorpusRefusedRuns}{0}
\newcommand{\CorpusDeopted}{10}
\newcommand{\CorpusFullAdmit}{1347}
\newcommand{\CorpusFullSame}{1347}
\newcommand{\CorpusFullUnsound}{0}
\newcommand{\CorpusFullRefusedRuns}{10}
\newcommand{\HousesAll}{461}
\newcommand{\HousesTest}{461}
\newcommand{\HousesLifted}{168}
\newcommand{\HousesCompiled}{168}
\newcommand{\HousesStatements}{208}
\newcommand{\HousesChanges}{377}
\newcommand{\HousesComplete}{48}
\newcommand{\HousesRuns}{1848}
\newcommand{\HousesAdmit}{1841}
\newcommand{\HousesAdmitPct}{100}
\newcommand{\HousesSame}{1841}
\newcommand{\HousesUnsound}{0}
\newcommand{\HousesStopped}{0}
\newcommand{\HousesStoppedRuns}{0}
\newcommand{\HousesRefusedCrash}{7}
\newcommand{\HousesRefusedRuns}{0}
\newcommand{\HousesDeopted}{3}
\newcommand{\HousesFullAdmit}{1838}
\newcommand{\HousesFullSame}{1838}
\newcommand{\HousesFullUnsound}{0}
\newcommand{\HousesFullRefusedRuns}{3}
\newcommand{\CorpusBothScripts}{733}
\newcommand{\CorpusBothRuns}{3212}
\newcommand{\CorpusBothSame}{3198}
\newcommand{\CorpusBothLifted}{292}
\newcommand{\ComposeFacts}{38}
\newcommand{\ComposeGuaranteed}{33}
\newcommand{\ComposeStageOneStops}{9}
\newcommand{\ComposeLikeFresh}{196}
\newcommand{\ComposeMonths}{196}
\newcommand{\ComposeBroken}{0}
\newcommand{\OhMonths}{92}
\newcommand{\OhTrusted}{0.88}
\newcommand{\OhFull}{0.88}
\newcommand{\OhDeopt}{0.88}
\newcommand{\OhPct}{$-1$}
\newcommand{\OhPctLow}{$-3$}
\newcommand{\OhPctHigh}{$+1$}
\newcommand{\OhCompile}{0.03}
\newcommand{\OhCompileMs}{28}
\newcommand{\OhFresh}{0.94}
\newcommand{\ImgBatches}{30}
\newcommand{\ImgBFpCorrect}{5}
\newcommand{\ImgBFpNeedless}{25}
\newcommand{\ImgBFpUnsound}{0}
\newcommand{\ImgBFpCorrectPct}{17}
\newcommand{\ImgBFpNeedlessPct}{83}
\newcommand{\ImgBFullCorrect}{10}
\newcommand{\ImgBFullNeedless}{20}
\newcommand{\ImgBFullUnsound}{0}
\newcommand{\ImgBFullCorrectPct}{33}
\newcommand{\ImgBFullNeedlessPct}{67}
\newcommand{\ImgBDeoptCorrect}{25}
\newcommand{\ImgBDeoptNeedless}{5}
\newcommand{\ImgBDeoptUnsound}{0}
\newcommand{\ImgBDeoptCorrectPct}{83}
\newcommand{\ImgBDeoptNeedlessPct}{17}
\newcommand{\ImgBTrustedUnsound}{15}
\newcommand{\ImgBDeopted}{15}
\newcommand{\CcPrograms}{11}
\newcommand{\CcAnalysesMin}{14}
\newcommand{\CcAnalysesMax}{207}
\newcommand{\CcFullAnalysesMax}{122}
\newcommand{\CcMsMin}{21}
\newcommand{\CcMsMax}{35}
\newcommand{\CcRatioMax}{2.3}
\newcommand{\CsPrograms}{8}
\newcommand{\CsFullMin}{6}
\newcommand{\CsFullMax}{17}
\newcommand{\CsItems}{23}
\newcommand{\CsItemMin}{1}
\newcommand{\CsItemMax}{2}
\newcommand{\CsAdmissionTypesOnly}{yes}
\newcommand{\CsTrafficCols}{6}
\newcommand{\CsTrafficFileCols}{9}
\newcommand{\CsBikesCols}{7}
\newcommand{\CsBikesFileCols}{16}
\newcommand{\CsTaxiCols}{9}
\newcommand{\CsTaxiFileCols}{20}
\newcommand{\RandFull}{127}
\newcommand{\RandDeopt}{361}
\newcommand{\RandDeopted}{234}
\newcommand{\RandFolders}{240}
\newcommand{\TestCount}{223}
\newcommand{\LeanLines}{1077}
\newcommand{\LeanDeoptLines}{178}
\newcommand{\LeanCoreLines}{899}
\newcommand{\SrcLines}{8,082}


\lstdefinelanguage{PrepLang}{
  morekeywords={LOAD,LOAD_TABLES,LOAD_IMAGE,LOAD_IMAGES,SAVE,SAVE_TABLES,SAVE_IMAGE,SAVE_IMAGES,
    AS,TO,WITH,ON,USING,WHERE,TYPE,EXCEPT,IF_EXISTS,RECURSIVE,DROP_NULL,FILL_NULL,NORMALIZE,
    FILTER,SELECT,DROP,EXPECT,RESIZE,GRAYSCALE,COMPRESS,CONVERT,REMOVE_CORRUPT,LINK,
    RANGE,MINMAX,ZSCORE,KEEP_ASPECT,QUALITY,JPG,PNG,NPY,NUMERIC,NOT_NULL,BETWEEN,UNIQUE,
    REFERENCES,SCHEMA_SAME_AS,CHANNELS,MIN_WIDTH,MIN_HEIGHT,MATCH_RATE,AND,OR,NOT},
  sensitive=false, comment=[l]{\#}, morestring=[b]"
}
\lstset{language=PrepLang, basicstyle=\ttfamily\footnotesize, keywordstyle=\bfseries,
  commentstyle=\itshape\color{gray}, numbers=left, numberstyle=\tiny, xleftmargin=1.5em,
  frame=tb, columns=fullflexible, breaklines=true}

\begin{document}

\title{PrepLang: Data-Aware Compilation of Tabular and\\Image Preprocessing Pipelines}

\author{\IEEEauthorblockN{Nikhil V}
\IEEEauthorblockA{School of Computer Science and Engineering\\
Vellore Institute of Technology, India\\
atnikhil06@gmail.com}}

\maketitle

\begin{abstract}
Data-preparation scripts for machine learning fail late or not at all: a misspelled column, a
stray word in a numeric column, a broken image file or floating-point pixels saved as an 8-bit
image surface only while the data is being processed, or never, as a silently corrupted dataset.
We present PrepLang, a small language and C++ compiler that treats the input data as a
compile-time input. The compiler derives an \emph{abstract dataset state}---column types, missing
values and value intervals, per-image sizes, channels, pixel type and format, and image--label
match rates---from the data itself (every CSV cell) or from metadata alone (Parquet footer
statistics, image headers), carries it through the pipeline with transfer functions, and uses it
to type-check every operation, to decide every declared data contract as proven, violated or
deferred to a run-time check, and to optimize. Proofs are bound to the inputs in two ways: a
fingerprint guard accepts the exact files, and a \emph{data certificate}---the facts the proofs
need, found by analysing the program again under weaker facts---lets the compiled program accept
new batches of data without recompiling. We state both soundness arguments. On \RbTotal{} data-preparation failures from GitHub issues, selected by error message,
PrepLang reports \RbEarly{} (\RbEarlyPct\%) before any data is processed; all \RbCases{} of our
reproductions behave as classified. On a 16-million-row Parquet file, compilation takes
\PqSixteen\,s whatever the file size, and reports injected mistakes in \ScTabNamePl\,s where the
equivalent script fails after \ScTabNameScript\,s; on 3,000 images a float-pixel mistake that the
script turns into 3,000 black images is reported in \ScImgFloatPl\,s. On 24 injected mistakes,
PrepLang reports all 24 before processing, while a script-style pipeline crashes part-way on 20 and
silently writes wrong output on 4. Of \CertBatches{} new batches for seven workloads, the
certificate admits \CertAdmitPct\% where the compiled file's own facts would admit
\FactsAdmitPct\%, and no admitted run differs from a fresh compilation. A parallel scanner analyses
one million CSV rows in \CsvMillion\,s, \CsvSpeedup$\times$ faster than pandas' default reader.
\end{abstract}

\begin{IEEEkeywords}
compilers, domain-specific languages, abstract interpretation, hybrid type checking, data
validation, data preprocessing, machine learning pipelines
\end{IEEEkeywords}

\section{Introduction}

Data preparation is a large part of building machine-learning systems, and its errors propagate
silently into models~\cite{sculley2015,polyzotis2018}. A preparation pipeline is usually a Python
script built on pandas~\cite{mckinney2010} and OpenCV~\cite{bradski2000}. Such a script cannot
know, until it runs, whether the column it normalizes holds numbers, whether a folder contains a
grayscale image among color ones, or whether every image has a label. When an assumption is wrong,
the script raises an exception after processing part of the data or, worse, finishes and writes a
wrong dataset: a threshold applied to an already normalized column removes every row, and a float
image in $[0,1]$ saved as PNG comes out black.

Data validation tools~\cite{breck2019,schelter2018,bantilan2020,gx} check declared expectations on
the data, but at run time and on a dataset rather than on the program that transforms it. Static
analyzers for data-science code~\cite{urban2019,subotic2022,lagouvardos2020,zhuang2022} look at the
program but mostly not at the data. An analysis of R scripts~\cite{gerstl2026} reads the datasets
but infers only the shape of data frames, and compile-time data contracts~\cite{mirji2026} check
the structure of declared types. None of them knows whether a column has missing values, what
range its values take, or how large an image is.

Batch preparation has a property these approaches leave unused: \emph{the inputs exist when the
pipeline is compiled}. PrepLang makes the data a compile-time input, as F\# type providers do for
types~\cite{petricek2016}, and contributes:

\begin{itemize}
\item \textbf{One abstract state for tables and images}, computed from the data or from metadata
only---Parquet footer statistics and image headers, at a cost independent of the data size---and
propagated by transfer functions (Sections~\ref{sec:state} and~\ref{sec:metadata}).
\item \textbf{Hybrid checking for data contracts.} In the spirit of hybrid type
checking~\cite{flanagan2006}, each \texttt{EXPECT} is proven, violated (naming the offending rows
or files) or deferred to a run-time check. An exactness flag makes refutation sound, and a
fingerprint guard binds proofs to the exact inputs, so the checks of proven contracts can be
deleted. We give the soundness argument (Section~\ref{sec:contracts}).
\item \textbf{Data certificates.} The program is analysed again with weaker facts about its
inputs; the facts every proof survives without are dropped. The rest is a certificate, with the
proofs that need each fact, that the generated program checks to accept a new input. Monotone
transfer functions make this sound (Section~\ref{sec:cert}).
\item \textbf{State-driven optimization with an explicit exact/approximate boundary}
(Section~\ref{sec:opt}).
\item \textbf{An evaluation} with injected mistakes, performance and scale experiments, and a study
of \RbTotal{} real failures from GitHub issues with checked reproductions (Section~\ref{sec:eval}).
\end{itemize}

The compiler has about 8,500 lines of C++17 and 181 unit tests; the artifact includes all
evaluation scripts.

\section{A Motivating Example}\label{sec:example}

Fig.~\ref{fig:linked} prepares images and their labels for training. The folder holds six images of
different sizes, one of them grayscale, and one file (\texttt{broken.jpg}) that is not an image.

\begin{figure}[t]
\begin{lstlisting}
LOAD_IMAGES "dataset/" AS imgs RECURSIVE;
LOAD "labels.csv" AS labels;
REMOVE_CORRUPT imgs;
LINK imgs WITH labels ON filename AS training;
EXPECT training MATCH_RATE 100%;
RESIZE imgs TO 128 128;
NORMALIZE imgs RANGE 0 1;   # pixels become FLOAT32
CONVERT imgs TO NPY;
SAVE_IMAGES imgs TO "processed/ml_ready/";
SAVE labels AS "processed/labels.csv";
\end{lstlisting}
\caption{Images linked to labels (\texttt{examples/images/linked.prep}).}
\label{fig:linked}
\end{figure}

Three small edits turn this into a typical broken pipeline. Without line~3, the equivalent script
fails on the broken file after decoding the others; PrepLang reports at compile time that
``1 of 7 files cannot be decoded''. With a labels file that lacks one image, the script's
assertion fails only after every image has been decoded; PrepLang reports ``match rate of training
is 83.33\%, expected at least 100\%''. With \texttt{CONVERT imgs TO PNG} in line~8, the script
finishes without any error and writes black images. PrepLang instead reports:

\begin{lstlisting}[language={},numbers=none]
error[type]: cannot save FLOAT32 pixels as PNG:
  NORMALIZE at line 7 scaled them to [0, 1]
  = hint: image files store 8-bit pixels, so these
    would come out almost black; add CONVERT imgs TO NPY
\end{lstlisting}

To produce this message, the compiler must know the pixel type and value range at line~9, which
come from the file headers read at line~1 and the effect of line~7.

\section{The PrepLang Language}\label{sec:lang}

PrepLang programs are straight-line sequences of statements over five kinds of data: a
\texttt{TABLE} (one CSV or Parquet file), a \texttt{TABLESET} (a folder of them), an
\texttt{IMAGE}, an \texttt{IMAGESET} (a folder, optionally recursive) and a \texttt{DATASET} (images
linked to a table by a file-name column). Table statements are \texttt{FILL\_NULL},
\texttt{DROP\_NULL}, \texttt{NORMALIZE}, \texttt{FILTER} (with \texttt{AND}, \texttt{OR},
\texttt{NOT} and comparisons), \texttt{SELECT} and \texttt{DROP}; image statements are
\texttt{RESIZE} (optionally \texttt{KEEP\_ASPECT}), \texttt{GRAYSCALE}, \texttt{NORMALIZE RANGE},
\texttt{CONVERT}, \texttt{COMPRESS} and \texttt{REMOVE\_CORRUPT}.

\emph{Selectors} address many columns at once: \texttt{company.*.salary} names the column in every
table of a table set, \texttt{company.*.* WHERE TYPE NUMERIC EXCEPT [id]} selects by the inferred
type, and \texttt{IF\_EXISTS} skips tables without the column. Selectors are resolved against the
inferred schemas, so a selector that matches a text column or misses a table is reported once, with
a note for every table it affects.

\emph{Contracts} are written with \texttt{EXPECT}: \texttt{NUMBER}, \texttt{STRING},
\texttt{BOOLEAN}, \texttt{NOT\_NULL}, \texttt{BETWEEN}, \texttt{UNIQUE} for columns;
\texttt{REFERENCES} and \texttt{SCHEMA\_SAME\_AS} across tables; \texttt{CHANNELS},
\texttt{MIN\_WIDTH}, \texttt{MIN\_HEIGHT} for images; and \texttt{MATCH\_RATE} for linked datasets.

Every statement begins with its own keyword, so the grammar is LL(1). The front end is a
hand-written lexer and a recursive-descent parser with panic-mode recovery; every diagnostic has a
caret, a hint and, for misspelled names, an edit-distance suggestion.

\section{The Abstract Dataset State}\label{sec:state}

The semantic analyzer maintains a state $\sigma$ mapping every dataset to an abstract value.
Programs have no loops, so the analysis is a single forward pass in the style of abstract
interpretation~\cite{cousot1977} and needs no widening.

\textbf{Columns.} For a table $T$ and column $c$, let $V_c(T)$ be the multiset of present values
and $M_c(T)$ the number of missing ones. A column is abstracted as
$a = \langle \tau, \nu, [l,u], \epsilon \rangle$: a type
$\tau \in \{\textsc{Integer}, \textsc{Float}, \textsc{String}, \textsc{Boolean}\}$; $\nu \in
\{\textit{none}, \textit{maybe}\}$ for missing values; an interval~\cite{cousot1976}; and an
exactness flag $\epsilon$. Its concretization is
\begin{multline*}
\gamma(a) = \{\, T \mid V_c(T) \subseteq [l,u] \text{ with type } \tau,\;
  \nu = \textit{none} \Rightarrow M_c(T) = 0, \\
  \epsilon \Rightarrow \big(M_c(T) > 0 \Leftrightarrow \nu = \textit{maybe}\big) \wedge
  \big(V_c(T) = \emptyset \vee \min V_c = l \wedge \max V_c = u\big) \,\}.
\end{multline*}
So an exact column promises two things: both interval bounds occur in the data, and the
missing-value flag is certain. For a CSV file, \texttt{LOAD} reads every cell and builds the
exact state: the most specific type consistent with all cells (a numeric-looking column that
becomes \textsc{String} remembers the offending line), the true minimum and maximum, and the true
number of missing values. Files that are not UTF-8, which pandas would reject, are reported with
their line and byte.

\textbf{Images.} An image set is abstracted as $\langle S, \delta, [l,u], \phi, q\rangle$: a list
$S$ of per-file shapes (width, height, channels, broken or not); the pixel type
$\delta \in \{\textsc{Uint8}, \textsc{Uint16}, \textsc{Float32}\}$; the pixel value range; the
output format $\phi \in \{\textsc{Source}, \textsc{Jpg}, \textsc{Png}, \textsc{Npy}\}$; and the
JPEG quality $q$. Shapes come from the PNG \texttt{IHDR} chunk or the JPEG start-of-frame marker;
the format is decided by the file's first bytes, as the decoder does, not by its name. A file whose
header cannot be parsed is broken.

\textbf{Links.} \texttt{LINK} records the numbers of images, labels and matches; the match rate is
$m/\max(|\mathit{images}|, |\mathit{labels}|)$.

\textbf{Transfer functions.} Table~\ref{tab:transfer} lists representative transfer functions.
\texttt{FILTER} refines intervals with the comparison and, following SQL's three-valued
logic~\cite{codd1979}, makes compared columns non-missing, because a comparison with a missing
value is unknown and the row is dropped; the generated code uses pandas' nullable types so that
execution follows the same rule. An empty refined interval means the filter removes every row,
which is reported. Using a column after \texttt{DROP} or \texttt{SELECT} removed it is reported
like a use after free. Save statements check that the format can hold the pixels:
\textsc{Float32} pixels cannot be saved as JPG or PNG, and 4-channel or \textsc{Uint16} images
cannot become JPG.

\begin{table}[t]
\caption{Representative transfer functions}
\label{tab:transfer}
\centering
\footnotesize
\begin{tabular}{@{}ll@{}}
\toprule
Statement & Effect on the abstract state \\
\midrule
\texttt{FILL\_NULL c WITH v} & $\nu \gets \textit{none}$; if $\nu$ was \textit{maybe}:
  $[l,u] \gets [l,u] \sqcup [v,v]$ \\
\texttt{DROP\_NULL t} & $\nu \gets \textit{none}$ for all columns; $\epsilon \gets$ false \\
\texttt{NORMALIZE c MINMAX} & $\tau \gets \textsc{Float}$; $[l,u] \gets [0,1]$; $\epsilon \gets$ false \\
\texttt{FILTER t WHERE c $\ge$ k} & $[l,u] \gets [\max(l,k), u]$; $\nu_c \gets \textit{none}$;
  $\epsilon \gets$ false \\
\texttt{RESIZE i TO w h} & every shape $\gets w \times h$ \\
\texttt{GRAYSCALE i} & every channel count $\gets 1$ \\
\texttt{NORMALIZE i RANGE a b} & $\delta \gets \textsc{Float32}$; $[l,u] \gets [a,b]$ \\
\texttt{CONVERT i TO f} & $\phi \gets f$ \\
\texttt{REMOVE\_CORRUPT i} & remove broken files from $S$ \\
\bottomrule
\end{tabular}
\end{table}

\section{Metadata-Only Analysis}\label{sec:metadata}

Reading every cell makes compilation cost as much as reading the data once
(Section~\ref{sec:eval}). For Parquet files~\cite{parquet,melnik2010} the compiler reads only the
footer, a Thrift-encoded structure at the end of the file, and builds the state from the schema,
the row count and the per-row-group column statistics, which databases already use to skip data
as \emph{zone maps}~\cite{moerkotte1998}:
\begin{itemize}
\item $\tau$ from the physical and logical types (dates, decimals and nested columns are
reported as unsupported);
\item $[l,u]$ from the smallest \texttt{min\_value} and largest \texttt{max\_value} over row groups,
decoded with the column's signedness. The format only promises \emph{bounds}; the interval is exact
only when every row group sets \texttt{is\_min\_value\_exact} and \texttt{is\_max\_value\_exact},
which current writers such as Arrow do for numbers;
\item $\nu$ from the summed \texttt{null\_count}. pandas reads a floating-point NaN as a missing
value, but the null count excludes NaN, so a float column is \textit{none} only if the file also
records a NaN count;
\item $\epsilon$ only if the interval is exact and the number of missing values is known.
\end{itemize}
Contracts that need the values themselves (\texttt{UNIQUE}, \texttt{REFERENCES}, \texttt{LINK}) are
deferred. The fingerprint of a Parquet input is the CRC-32 of its footer: the proofs use nothing
else. The cost is independent of the number of rows. Soundness now rests on one assumption
(\emph{S}): the writer's statistics are truthful.

\section{Hybrid Contract Checking and Soundness}\label{sec:contracts}

For a contract $\varphi$ at a program point with state $\sigma$, the compiler returns
\textsc{Proven} if every dataset in $\gamma(\sigma)$ satisfies $\varphi$ (no check is generated),
\textsc{Violated} if none does (compilation fails), and \textsc{Deferred} otherwise (a run-time
check is generated). This is hybrid type checking~\cite{flanagan2006} and soft
typing~\cite{cartwright1991} applied to data: prove what the state decides, check the rest. For
\texttt{BETWEEN a AND b}, the contract is proven if $[l,u] \subseteq [a,b]$ and violated if the
column is exact and $[l,u] \not\subseteq [a,b]$; for a CSV input the compiler then names the first
offending row. A deferred contract is an assumption for the rest of the program, because execution
continues only if the check passed: $\sigma$ is refined with $\varphi$ (for \texttt{BETWEEN},
$[l,u] \gets [l,u] \sqcap [a,b]$ and $\epsilon \gets$ false, since $a$ and $b$ need not occur).

\textbf{Fingerprint guard.} The generated program recomputes a CRC-32 of every CSV input, of every
Parquet footer and of every image folder listing (names and sizes), and stops if any differs from
the value recorded at compile time.

\begin{lemma}[Initial state]\label{lem:init}
If $T$ is the content of a CSV input, $T \in \gamma(\sigma_0)$. If $T$ is the content of a
Parquet input whose statistics are truthful (S), $T \in \gamma(\sigma_0)$.
\end{lemma}
\begin{IEEEproof}
For CSV, $\sigma_0$ is computed from all cells, so every clause holds with equality. For Parquet,
S gives $V_c \subseteq [l,u]$; exactness is claimed only when the attained-bound flags are set and
missing values are fully counted, so both clauses guarded by $\epsilon$ hold.
\end{IEEEproof}

\begin{theorem}[Transfer soundness]\label{thm:transfer}
For every statement $s$ with semantics $\llbracket s \rrbracket$ and transfer function $f_s$,
$T \in \gamma(\sigma) \Rightarrow \llbracket s \rrbracket(T) \in \gamma(f_s(\sigma))$.
\end{theorem}
\begin{IEEEproof}[Proof sketch]
By cases; we show the ones that touch exactness. \texttt{FILTER} keeps a subset of rows, so
$V'_c \subseteq V_c \subseteq [l,u]$, and every kept row satisfies the condition, which gives the
refined bound; a compared column has no missing values in kept rows (three-valued logic); bounds
may no longer occur, so $\epsilon$ is cleared. \texttt{FILL\_NULL c WITH v} gives
$V'_c = V_c \cup \{v\}$ if $M_c > 0$ and $V'_c = V_c$ otherwise; the joined interval covers both.
If $\epsilon$ holds, $M_c > 0 \Leftrightarrow \nu = \textit{maybe}$, so $v$ is added exactly when
the transfer function adds it, and the bounds of the join occur. \texttt{NORMALIZE},
\texttt{DROP\_NULL} and the image statements change values or shapes in a way the table states
exactly, or clear $\epsilon$.
\end{IEEEproof}

\begin{theorem}[Verdict soundness]\label{thm:verdict}
If the verdict for $\varphi$ is \textsc{Proven}, every $T \in \gamma(\sigma)$ satisfies $\varphi$;
if it is \textsc{Violated}, no $T \in \gamma(\sigma)$ satisfies $\varphi$.
\end{theorem}
\begin{IEEEproof}[Proof sketch]
For \texttt{BETWEEN}: $[l,u] \subseteq [a,b]$ and $V_c \subseteq [l,u]$ give $V_c \subseteq
[a,b]$. If the column is exact and, say, $u > b$, then $u \in V_c$ and $\varphi$ fails. For
\texttt{NOT\_NULL}: \textit{none} gives $M_c = 0$; a violation is reported only for an exact column
with a positive missing count from the file. Image contracts are decided per file on exact shapes.
\end{IEEEproof}

\begin{corollary}[Check elimination]\label{cor:checks}
Barring a CRC-32 collision, every run that passes the fingerprint guard reaches each
\texttt{EXPECT} with a dataset in $\gamma(\sigma)$, so deleting the checks of proven contracts does
not change which runs succeed.
\end{corollary}
\begin{IEEEproof}
The guard makes the input equal to the analyzed one (for Parquet: its footer, which with S is all
the proofs use); apply Lemma~\ref{lem:init}, Theorem~\ref{thm:transfer} along the program (deferred
contracts refine $\sigma$ only after their check passed) and Theorem~\ref{thm:verdict}.
\end{IEEEproof}

Writing down $\gamma$ was useful in practice: it exposed an unsound case in our first Parquet
implementation, which marked a float column exact although its missing values were not known, so
that \texttt{FILL\_NULL} followed by \texttt{BETWEEN} could be reported as violated when the fill
value never occurs. The case is now excluded by the definition of $\epsilon$ and covered by a unit
test. CRC-32 detects accidental changes, not deliberate ones; a cryptographic hash can replace it.

\section{Data Certificates}\label{sec:cert}

The fingerprint guard is sound but strict: a preparation pipeline usually runs again on tomorrow's
file, which has the same layout and different values, and the guard refuses it. Most facts of the
compiled file are not needed by any proof, though. In the running example the data has marks in
$[35,95]$, but the only proof about marks is \texttt{BETWEEN 0 AND 100}, which needs $[0,100]$.

\textbf{Certificate.} A certificate is an input state $\sigma_C$ in which every column is inexact
($\epsilon$ false, Section~\ref{sec:state}): a set of tables, not one table. An input $T$ is
\emph{admitted} if $T \in \gamma(\sigma_C)$: every column read keeps its type (\textsc{Integer}
values may fill a \textsc{Float} column, and a column without values fits any type), has no missing
values where $\sigma_C$ says \textit{none}, and stays inside the certified interval. Let
$\Pi(\sigma)$ be the proven contracts and removed operations the analysis derives from input state
$\sigma$, leaving out contracts proven by reading the values (\texttt{UNIQUE},
\texttt{REFERENCES}, \texttt{MATCH\_RATE}).

\begin{lemma}[Monotonicity]\label{lem:mono}
Order states pointwise ($\textit{none} \sqsubseteq \textit{maybe}$, interval inclusion, $\epsilon$
false). For every statement, $\sigma' \sqsubseteq \sigma$ implies $f_s(\sigma') \sqsubseteq
f_s(\sigma)$, and every proof and removal decided in $\sigma$ is also decided in $\sigma'$.
\end{lemma}
\begin{IEEEproof}[Proof sketch]
The decisions are inclusions of the state in a fixed predicate ($[l,u] \subseteq [a,b]$,
$\nu = \textit{none}$, a comparison true on all of $[l,u]$ without missing values), which a smaller
state keeps. Transfer functions are joins and meets with constants, except where a decision changes
the path: a \texttt{FILTER} that is always true in $\sigma'$ leaves $\sigma'$ unchanged, and
$\sigma' \sqsubseteq \sigma \sqcap \textit{cond} = f_s(\sigma)$ because $\sigma'$ already satisfies
the condition; a contract proven in $\sigma'$ but deferred in $\sigma$ leaves $\sigma'$, and
$\sigma' \sqsubseteq \sigma \sqcap \varphi$ likewise. \texttt{FILL\_NULL} adds $v$ only when
$\nu = \textit{maybe}$, which is monotone because $[l,u] \subseteq [l,u] \sqcup [v,v]$.
\end{IEEEproof}

\begin{theorem}[Certificate soundness]\label{thm:cert}
If $\Pi(\sigma_C) \supseteq \Pi(\sigma_0)$ and an input $T$ is admitted, every removal and every
proof the compiled program relies on is valid for $T$.
\end{theorem}
\begin{IEEEproof}
Let $\alpha(T)$ be the inexact state of $T$; admission means $\alpha(T) \sqsubseteq \sigma_C$. By
Lemma~\ref{lem:mono}, $\Pi(\alpha(T)) \supseteq \Pi(\sigma_C) \supseteq \Pi(\sigma_0)$, and by
Lemma~\ref{lem:init}, Theorem~\ref{thm:transfer} and Theorem~\ref{thm:verdict} applied to
$\alpha(T)$, each of these proofs holds for $T$ and each removed operation does nothing on $T$.
Checks of contracts proven from the values are kept in the generated program and run whenever an
input differs from the compiled one.
\end{IEEEproof}

\textbf{Finding the certificate.} The compiler starts from $\sigma_0$ with every $\epsilon$
cleared and first checks that this still gives $\Pi(\sigma_0)$; otherwise only the exact files are
accepted. Then, for every column the optimized program reads, it tries to weaken one fact at a time
and analyses the whole program again: missing values are allowed; a lower bound is replaced by
$-\infty$, or else by the loosest program constant below it that works; likewise for the upper
bound. A weakening is kept when $\Pi$ is unchanged. The analysis compares intervals only with the
program's constants (contract bounds, comparison values, fill values), so no other bound can keep
more proofs; the result is 1-minimal: no remaining fact can be weakened on its own. A rejected
weakening tells which proofs need the fact, and the certificate lists them, so a refused input is
explained by the row and the proof:

\begin{lstlisting}
marks is 104 in data row 4, above 100; the compiled
program relies on this for line 5: EXPECT
students.marks BETWEEN 0 AND 100
\end{lstlisting}

A column costs at most $1 + 2(k+1)$ analyses for $k$ constants. The analyses read no data: tables
are replaced by their facts, and image headers are cached. Image folders keep the exact fingerprint.

\section{State-Driven Optimization}\label{sec:opt}

Analysis results are attached to a three-address intermediate representation in which wildcards
are already expanded per table. Three passes and the code generator use them.
\textbf{No-effect elimination} removes operations the state proves useless: \texttt{FILL\_NULL} on
a column without missing values, a \texttt{FILTER} that is always true, \texttt{GRAYSCALE} on
1-channel images, \texttt{CONVERT} to the current format.
\textbf{Check elimination} removes checks of proven contracts (Corollary~\ref{cor:checks}); the
check of a contract proven from the values stays, guarded so that it runs only when an input is not
the compiled file (Section~\ref{sec:cert}).
\textbf{Dead-code elimination} removes unused temporaries and image operations after the last save.
\textbf{Dead-store elimination and column pruning} run a backward liveness analysis over columns;
\texttt{LOAD} reads only live columns, and \texttt{DROP\_NULL} counts as reading every column,
since pruning a column before it would change which rows are removed.
\textbf{Image streaming} turns all image operations up to a save into one loop that decodes,
transforms and writes one image at a time; this is exact because the operations act on each image
independently.

All of these preserve the output exactly. \textbf{RESIZE fusion}, which replaces two consecutive
resizes by one, does not: interpolation makes $\mathit{resize}_{224} \circ \mathit{resize}_{512}$
differ slightly from $\mathit{resize}_{224}$. PrepLang reports it as a suggestion and applies it
only with \texttt{--approx}: a preprocessing compiler should not change training data silently.

\section{Implementation}

The compiler is written in C++17 without external libraries: lexer, parser, AST printer and
Graphviz export, schema provider (CSV inference, a Thrift compact-protocol reader for Parquet
footers, PNG/JPEG header reading, CRC-32), semantic analyzer with selector resolver and contract
analyzer, lineage tracker, IR builder, pass manager, and an execution planner that merges the
pandas and OpenCV code generators in program order. Flags print every phase; \texttt{--explain}
prints the execution plan and \texttt{--dry-run} the files a run would read and write, without
running anything. The generated program writes a manifest with inputs, fingerprints, contract
verdicts, operations and outputs.

The CSV reader reads a file once and cuts it at line breaks into one piece per core. Each thread
splits its lines in place (\texttt{std::string\_view}), infers cell types and parses numbers in one
pass, using Clinger's exact fast path~\cite{clinger1990} when the digits and the power of ten are
exact doubles; the CRC-32 (slicing-by-8) runs on its own thread. Pieces are merged in file order,
keeping the first problem and the first cell that makes a column \textsc{String}, so the result
equals a sequential pass; a test compares both on files cut into 1 to 16 pieces, and a
differential test against the earlier sequential reader found no difference on 60 random files of
up to 250,000 rows with injected errors.

\section{Evaluation}\label{sec:eval}

We ask seven questions. \textbf{RQ1:} Does compiling against the data find injected mistakes before
any data is processed? \textbf{RQ2:} How many contracts are decided at compile time?
\textbf{RQ3:} What do the optimizations gain at run time? \textbf{RQ4:} What does analyzing the
data at compile time cost? \textbf{RQ5:} How early are mistakes reported on large inputs?
\textbf{RQ6:} How many real data-preparation failures would PrepLang report?
\textbf{RQ7:} Can a compiled program be reused on new data?

\textbf{Setup.} Intel Core i5-12450H (8 cores), 16\,GB RAM, Windows~11; MSVC (toolset 14.44,
Release); Python 3.14, pandas 3.0, NumPy 2.5, OpenCV 5.0, pyarrow 25. Times are medians of three
runs and include starting Python, as a user experiences them.

\textbf{Baselines.} A \emph{script-style interpreter} executes PrepLang programs statement by
statement with pandas and OpenCV, with default type inference, images held in memory, and each
\texttt{EXPECT} as an assertion. For image pipelines we add a \emph{hand-written streaming script}
(\texttt{eval/handwritten.py}) that processes one image at a time, as a careful programmer would.

\subsection{RQ1: Injected mistakes}

We injected 24 mistakes into the five example programs, one per mutant, in the spirit of mutation
testing~\cite{jia2011}: \emph{name} (misspelled columns), \emph{type} (normalizing text, comparing a
number with text, text into an integer column, zero-size resize), \emph{lifecycle} (use after
removal, a table operation on images), \emph{contract} (range, not-null, foreign key, schema,
channels, width, match rate), \emph{validity} (a broken image) and \emph{value} (well-typed programs
with silently wrong output, detected by oracles). The unmodified programs are controls.

\begin{table}[t]
\caption{When each injected mistake is found (24 mutants)}
\label{tab:rq1}
\centering
\footnotesize
\begin{tabular}{@{}lrrr@{}}
\toprule
& & \multicolumn{2}{c}{Script-style baseline} \\
\cmidrule(l){3-4}
Class & Mutants & Crash while running & Silent wrong output \\
\midrule
name       & 3 & 3 & 0 \\
type       & 5 & 4 & 1 \\
lifecycle  & 3 & 3 & 0 \\
contract   & 9 & 9 & 0 \\
validity   & 1 & 1 & 0 \\
value      & 3 & 0 & 3 \\
\midrule
total      & 24 & 20 & 4 \\
\bottomrule
\end{tabular}
\end{table}

PrepLang reports all 24 at compile time (23 errors, one warning) and nothing on the controls; the
baseline crashes part-way on 20 and completes 4 with wrong output (Table~\ref{tab:rq1}): text
filled into an integer column, a threshold on a normalized column that removes every row (PrepLang:
``FILTER removes every row: the condition is never true''), and two pipelines that write black
images. Compilation took \RqOneMsMin--\RqOneMsMax\,ms (median \RqOneMsMed\,ms). The mutants are our
own, so RQ1 shows which kinds of mistakes are caught; RQ6 measures how often real ones are.

\subsection{RQ2: Contracts decided at compile time}

The controls declare 16 contracts. PrepLang proves 15 (94\%) and defers one: in
\texttt{students.prep}, \texttt{EXPECT marks BETWEEN 60 AND 100} follows a filter that guarantees
only $[50, 95]$.

\subsection{RQ3: Run-time effect of the optimizations}

Table~\ref{tab:rq3} shows three pipelines: a 200,000-row, 40-column CSV (58\,MB) with two contracts,
a filter and a selection of four columns, where the optimizer proves both contracts and reads five
columns; 1,000 images resized to $224\times224$, made gray and saved as JPG; and the same images
resized twice and saved as PNG, with and without \texttt{--approx}.

\begin{table}[t]
\caption{Run-time effect of the optimizations}
\label{tab:rq3}
\centering
\footnotesize
\begin{tabular}{@{}lrr@{}}
\toprule
\multicolumn{3}{@{}l}{\emph{Table pipeline, 200k rows $\times$ 40 columns}} \\
Variant & Columns read, checks & Time (s) \\
\midrule
PrepLang, optimized & 5, 0 & \TabOpt \\
PrepLang, \texttt{--O0} & 40, 2 & \TabOz \\
Script-style baseline & 40, 2 & \TabBase \\
\midrule
\multicolumn{3}{@{}l}{\emph{Image pipeline, 1,000 images}} \\
Variant & Peak memory (MiB) & Time (s) \\
\midrule
PrepLang (streaming) & \ImgPlMem & \ImgPl \\
Hand-written streaming script & \ImgHwMem & \ImgHw \\
Script-style baseline (all in memory) & \ImgBaseMem & \ImgBase \\
\midrule
\multicolumn{3}{@{}l}{\emph{RESIZE fusion, 1,000 images}} \\
Variant & PSNR vs.\ exact (dB) & Time (s) \\
\midrule
Exact (default) & -- & \FusExact \\
\texttt{--approx} & \FusPsnr{} (min.\ \FusPsnrMin) & \FusApprox \\
\bottomrule
\end{tabular}
\end{table}

For the table, both variants write byte-identical output. Starting Python and importing pandas
takes \TabStart\,s of every run; without it, the optimized program processes the data in
\TabProcOpt\,s instead of \TabProcOz\,s (\TabCut\% less). For images, PrepLang's streaming matches
the memory of the careful hand-written script and is \ImgBaseMem/\ImgPlMem{} MiB below loading all
images; the benefit over a careful programmer is that streaming is automatic, not that it is
faster. RESIZE fusion saves \FusSave\% of the time, because decoding and encoding dominate, and
changes the output (PSNR~\cite{huynh2008} \FusPsnr\,dB): a small gain for a change to training data,
which supports keeping it opt-in. While building these benchmarks we found and fixed two
inefficiencies in the generated code: passing nullable types to \texttt{read\_csv} made reading six
times slower, and defensive copies after \texttt{FILTER} and \texttt{SELECT} cost up to 0.3\,s on
16 million rows.

\subsection{RQ4: Cost of the analysis}

Table~\ref{tab:rq4} shows compile time (with \texttt{--explain}, which analyzes but writes nothing)
next to one plain read of the same data in Python.

\begin{table}[t]
\caption{Compile time vs.\ one plain read of the data (s)}
\label{tab:rq4}
\centering
\footnotesize
\begin{tabular}{@{}lrr@{}}
\toprule
Input & Compile & Read in Python \\
\midrule
CSV, 10k rows $\times$ 10 columns & \CsvTenK & \CsvTenKRead \\
CSV, 100k rows $\times$ 10 columns & \CsvHundredK & \CsvHundredKRead \\
CSV, 1M rows $\times$ 10 columns & \CsvMillion & \CsvMillionRead \\
Parquet, 1M rows (\PqOneSize\,MiB) & \PqOne & \PqOneRead \\
Parquet, 4M rows (\PqFourSize\,MiB) & \PqFour & \PqFourRead \\
Parquet, 16M rows (\PqSixteenSize\,MiB) & \PqSixteen & \PqSixteenRead \\
100 JPG images, $640\times480$ & \ImgHundred & \ImgHundredRead \\
1,000 JPG images & \ImgThousand & \ImgThousandRead \\
3,000 JPG images & \ImgThreeK & \ImgThreeKRead \\
\bottomrule
\end{tabular}
\end{table}

With the parallel scanner, analysing every CSV cell of one million rows takes \CsvMillion\,s,
\CsvSpeedup$\times$ faster than pandas' default reader reads the same file and about as fast as
pandas with the multi-threaded pyarrow reader (\CsvArrowRead\,s, measured inside a running Python
process, so without process start); the earlier sequential reader took 2.04\,s. Metadata-only analysis goes further: Parquet compile time
stays at \PqSixteen\,s from one to 16 million rows, and reading image headers is \HdrRatio$\times$
faster than decoding 3,000 images.

\subsection{RQ5: Large inputs}

Table~\ref{tab:rq5} compares, for injected mistakes on large inputs, when PrepLang reports them
with when the equivalent script stops (or finishes with wrong output). The table pipeline fills,
normalizes, filters and selects 16 million Parquet rows; the image pipeline resizes, converts to
gray and saves 3,000 images, compared with the hand-written streaming script.

\begin{table}[t]
\caption{Time until a mistake is reported (s)}
\label{tab:rq5}
\centering
\footnotesize
\begin{tabular}{@{}lrr@{}}
\toprule
Case & PrepLang & Script \\
\midrule
16M rows: misspelled column in the last step & \ScTabNamePl & \ScTabNameScript \\
16M rows: range contract broken by the data & \ScTabContractPl & \ScTabContractScript \\
3,000 images: broken file, sorted last & \ScImgBrokenPl & \ScImgBrokenScript \\
3,000 images: float pixels saved as PNG & \ScImgFloatPl & \ScImgFloatScript$^{*}$ \\
\midrule
16M rows, correct: compile + run & \ScTabOkPl{} + \ScTabOkRun & \ScTabOkScript \\
3,000 images, correct: compile + run & \ScImgOkPl{} + \ScImgOkRun & \ScImgOkScript \\
\bottomrule
\multicolumn{3}{@{}l}{$^{*}$ finishes without an error; every written image is black}
\end{tabular}
\end{table}

Mistakes are reported in a fraction of a second, before any row or pixel is processed, while the
scripts fail only after the data has been read or every image processed. On correct programs,
PrepLang's generated code is somewhat slower than the scripts, mainly because pandas' nullable
types, which implement the missing-value semantics the analysis relies on, are slower than plain
NumPy columns.

\subsection{RQ6: Real data-preparation failures}

\textbf{Protocol.} We chose 17 error messages that Python libraries print for common
data-preparation failures (for example Pillow's \emph{cannot write mode RGBA as JPEG}, pandas'
\emph{not in index}, scikit-learn's \emph{feature names should match}); they include failures
PrepLang is not designed for (dates, text encodings, sample counts), so that the result measures
coverage. For each message we searched GitHub issues and read the first 12 results in order. An
issue qualified if the failure happened while loading, transforming or saving data files and its
root cause was stated; we kept at most three per message, giving \RbTotal{} issues (91 other
results were read and excluded, each with a recorded reason). Each qualifying issue was classified
as reported at compile time, reported once the requirement is declared with \texttt{EXPECT},
impossible in PrepLang by construction, or not covered; every positive classification was
reproduced as a PrepLang program with data mirroring the issue and checked against the compiler
(\RbPass{} of \RbCases{} reproductions behave as classified). This follows the practice of
bug-characteristic studies that mine issue trackers~\cite{islam2019,humbatova2020}.

\begin{table}[t]
\caption{Real failures from GitHub issues (\RbTotal{} issues)}
\label{tab:rq6}
\centering
\footnotesize
\begin{tabular}{@{}lr@{}}
\toprule
Outcome & Issues \\
\midrule
Reported at compile time, no annotation & \RbCompile \\
Reported at compile time with a declared \texttt{EXPECT} & \RbContract \\
Cannot occur in a PrepLang program & \RbDesign \\
Not covered & \RbNone \\
\quad input format (HEIC, PDF, audio, GeoTIFF, \ldots) & \RbFormat \\
\quad operation not in the language (groupby, pivot, \ldots) & \RbOperation \\
\quad dates & \RbDates \\
\quad one label file per image & \RbLabels \\
\quad input produced during the run & \RbProduced \\
\quad no suitable contract & \RbContractMissing \\
\quad model code & \RbModel \\
\bottomrule
\end{tabular}
\end{table}

\textbf{Results.} PrepLang reports \RbEarly{} of the \RbTotal{} failures (\RbEarlyPct\%) before any
data is processed (Table~\ref{tab:rq6}). The caught ones are the failures its state was designed
for: broken or misnamed image files, text in numeric columns, missing columns, gray or RGBA images
where a model needs RGB, keys missing from another table, images without labels, and files that are
not UTF-8. The study also changed the compiler: two classes were caught only after we decided image
formats by content (a transparent PNG screenshot saved as \texttt{.jpg} had been reported as a broken
file) and checked CSV encodings; before these fixes the rate was \RbBefore{} of \RbTotal{}
(\RbBeforePct\%), a number we measured by running the reproductions against the earlier compiler.
The uncovered failures point to extensions: more input formats, dates, per-image label files, and
set-valued operations such as pivots.

\subsection{RQ7: Reusing a compiled program on new data}

Seven workloads are compiled once: the four table examples, the same program twice on one file
(\texttt{optimize}), the Parquet example, the images linked to labels, and a sensor pipeline of
20,000 rows per day. Each input table then gets new batches: other days and another season of the
sensor process, and, for the examples, rows resampled with replacement, values moved by noise of
0.1, 0.5 and 1.5 standard deviations, 20\% missing values in one column, and one value far out of
range. Each batch is run by the compiled program (\texttt{table=batch.csv}) and, as ground truth,
compiled and run fresh. We compare three admission policies: the fingerprint, the compiled file's
own facts (type, no missing values, and $[\min,\max]$ of every column read, as a schema inferred by
TFDV~\cite{breck2019} would require), and the certificate.

\begin{table}[t]
\caption{New batches admitted by each policy (RQ7)}
\label{tab:rq7}
\centering
\footnotesize
\begin{tabular}{@{}lrrrr@{}}
\toprule
Workload & Batches & File facts & Certificate & Ends like fresh \\
\midrule
daily & 8 & 0 & 8 & 8 \\
students & 24 & 7 & 14 & 14 \\
optimize & 24 & 7 & 15 & 15 \\
ml\_split & 48 & 15 & 42 & 42 \\
tableset & 96 & 41 & 96 & 96 \\
parquet & 24 & 10 & 14 & 14 \\
linked\_images & 20 & 16 & 20 & 20 \\
\midrule
all & 244 & 96 & 209 & 209 \\
\bottomrule
\end{tabular}
\end{table}


The certificate admits \CertAdmit{} of \CertBatches{} batches (\CertAdmitPct\%), the file's own
facts \FactsAdmit{} (\FactsAdmitPct\%) and the fingerprint none. Every admitted run ends like the
fresh compilation (\CertAgree{} of \CertAdmit{}: the same output files, or both stopped by the
same contract), and no admitted run finished where the fresh compilation failed. Every refused
batch also fails to compile fresh (\CertFalse{} refused batches would have compiled), so on these
workloads the certificate is not only sound but loses nothing to its 1-minimal search. The sensor
workload shows the intended use: other days and a season 12 degrees warmer are all admitted, while
the file's own facts admit none of them.

During this experiment, the run-time \texttt{MATCH\_RATE} check, which had always been eliminated
before, ran for the first time on changed labels and failed where the compile-time proof held: it
counted a broken image that \texttt{REMOVE\_CORRUPT} had removed before the \texttt{LINK}. The
generated check now skips the same files the proof skipped, with a regression test.

\subsection{Threats to validity}

The error messages were chosen by the author, and one person classified the issues; the
reproductions confirm the classifications of positive cases but are constructed from the issue
text, not from the original data. GitHub's ranking favors some repositories, and 44 issues are few.
The mutants of RQ1 are our own, the baseline interpreter is ours, and performance numbers come from
one machine and synthetic data. Metadata-only soundness assumes truthful Parquet statistics. The
approach needs the inputs at compile time, which holds for batch preparation, and the language has
no joins, loops or user-defined functions. Certificates are 1-minimal, not minimum, and cover types,
missing values and intervals of tables; image inputs must stay exactly as compiled. The new batches
of RQ7 are generated edits of small example files and one synthetic process.

\section{Related Work}

\textbf{Data validation.} TensorFlow Data Validation~\cite{breck2019}, Deequ~\cite{schelter2018},
Great Expectations~\cite{gx}, Pandera~\cite{bantilan2020} and the Data Linter~\cite{hynes2017}
check expectations on datasets when a pipeline runs. PrepLang checks the same kinds of expectations
at a point in the program, against the state the preceding statements produce, and moves the check
to compile time when the state decides it.

\textbf{Static analysis of data-science code.} Urban~\cite{urban2019} argues for abstract
interpretation of data-science software, and Subotić et al.~\cite{subotic2022} built such an
analysis for notebooks. Zhuang and Lu~\cite{zhuang2022} type-check data-frame columns by abstract
interpretation, Lagouvardos et al.~\cite{lagouvardos2020} infer tensor shapes in TensorFlow
programs, and mlinspect~\cite{grafberger2022} inspects pipelines by instrumenting their execution.
Closest to our work, Gerstl et al.~\cite{gerstl2026} infer the shapes of R data frames and become
markedly more precise when the datasets are available; they infer column names and dimensions,
whereas PrepLang infers types, missing values and intervals for tables and per-file properties for
images, and uses them to decide contracts and to optimize. Sihler et al.~\cite{sihler2026} make
implicit assumptions of analysis scripts explicit so they can be checked at run time.

\textbf{Contracts and hybrid checking.} Soft typing~\cite{cartwright1991} and hybrid type
checking~\cite{flanagan2006} prove what they can statically and insert run-time checks for the rest;
PrepLang applies this to data contracts, where the facts come from the data. Mirji~\cite{mirji2026}
proves structural contracts for typed Spark pipelines at compile time and keeps a run-time schema
check at the output, since data may drift after compilation; PrepLang proves value-level contracts
from the data and handles drift by guarding the inputs, which is what allows deleting the checks.

\textbf{Reusing proofs on new inputs.} Proof-carrying code~\cite{necula1997} ships a proof with a
program so that a consumer can check it cheaply; PrepLang ships the facts its proofs depend on, and
the generated program checks the input against them. Dynamic compilers specialize code to observed
values behind guards and deoptimize when a guard fails~\cite{holzle1992}; progressive query
optimization attaches validity ranges to a plan and re-optimizes outside them~\cite{markl2004},
and parametric query optimization precomputes plans per parameter range~\cite{ioannidis1992}.
These guards protect performance and fall back to a slower path; a certificate protects proofs of
correctness and is derived by re-analysis rather than by cost estimates. Precondition
inference~\cite{cousot2013} computes conditions on a program's inputs that are necessary for its
assertions, and DeepInfer~\cite{ahmed2024} infers data preconditions of trained neural networks;
we infer sufficient conditions on input data for a compiled pipeline's proofs and removals, by
counterfactual re-analysis over the abstract state. TFDV~\cite{breck2019} validates new batches
against a schema inferred from earlier data; RQ7 shows that taking all facts of the data, rather
than the ones the program needs, refuses most new batches.

\textbf{Types from data, metadata and data-centric compilers.} F\# type
providers~\cite{petricek2016} infer types from sample data at compile time. Min/max statistics per
block are used by databases to skip data~\cite{moerkotte1998} and are stored in Parquet
footers~\cite{parquet}; PrepLang uses them to prove contracts. Halide~\cite{ragan2013} fuses
image-processing stages under a separate schedule, and Weld~\cite{palkar2017} optimizes analytics
across libraries through a common IR; their optimizations are driven by the program, whereas
PrepLang's are driven by facts about the data.

\section{Conclusion}

For batch data preparation, treating the input data as a compile-time input pays off: an abstract
dataset state computed from the data, or from its metadata alone, turns many data mistakes into
compile-time errors, decides almost all declared contracts statically, and justifies optimizations
that keep the output exact, with a fingerprint guard that keeps the proofs sound. Data
certificates carry those proofs to new batches: the compiled program accepts any input that has
the few facts its proofs use, and names the proof when it refuses one. On real failures
reported on GitHub, it catches about two in five before any data is processed; the rest outline
the next steps: more input formats, dates, per-image label files and set-valued operations.

\begin{thebibliography}{99}

\bibitem{sculley2015}
D. Sculley, G. Holt, D. Golovin, E. Davydov, T. Phillips, D. Ebner, V. Chaudhary, M. Young, J.-F.
Crespo, and D. Dennison, ``Hidden technical debt in machine learning systems,'' in \emph{Advances
in Neural Information Processing Systems (NeurIPS)}, 2015, pp. 2503--2511.

\bibitem{polyzotis2018}
N. Polyzotis, S. Roy, S. E. Whang, and M. Zinkevich, ``Data lifecycle challenges in production
machine learning: A survey,'' \emph{ACM SIGMOD Record}, vol. 47, no. 2, pp. 17--28, 2018.

\bibitem{mckinney2010}
W. McKinney, ``Data structures for statistical computing in Python,'' in \emph{Proc. 9th Python in
Science Conf. (SciPy)}, 2010, pp. 56--61.

\bibitem{bradski2000}
G. Bradski, ``The OpenCV library,'' \emph{Dr. Dobb's Journal of Software Tools}, 2000.

\bibitem{breck2019}
E. Breck, N. Polyzotis, S. Roy, S. E. Whang, and M. Zinkevich, ``Data validation for machine
learning,'' in \emph{Proc. Conf. on Machine Learning and Systems (MLSys)}, 2019.

\bibitem{schelter2018}
S. Schelter, D. Lange, P. Schmidt, M. Celikel, F. Biessmann, and A. Grafberger, ``Automating
large-scale data quality verification,'' \emph{Proc. VLDB Endowment}, vol. 11, no. 12,
pp. 1781--1794, 2018.

\bibitem{bantilan2020}
N. Bantilan, ``pandera: Statistical data validation of pandas dataframes,'' in \emph{Proc. 19th
Python in Science Conf. (SciPy)}, 2020, pp. 116--124.

\bibitem{gx}
Great Expectations, ``Great Expectations: data quality testing,'' \url{https://greatexpectations.io}.

\bibitem{hynes2017}
N. Hynes, D. Sculley, and M. Terry, ``The Data Linter: Lightweight, automated sanity checking for
ML data sets,'' in \emph{NIPS Workshop on Machine Learning Systems}, 2017.

\bibitem{urban2019}
C. Urban, ``Static analysis of data science software,'' in \emph{Static Analysis Symposium (SAS)},
LNCS 11822, 2019, pp. 17--23.

\bibitem{subotic2022}
P. Suboti\'{c}, L. Miliki\'{c}, and M. Stoji\'{c}, ``A static analysis framework for data science
notebooks,'' in \emph{Proc. 44th Int. Conf. on Software Engineering: Software Engineering in
Practice (ICSE-SEIP)}, 2022, pp. 13--22.

\bibitem{zhuang2022}
Y. Y. Zhuang and M.-Y. Lu, ``Enabling type checking on columns in data frame libraries by abstract
interpretation,'' \emph{IEEE Access}, vol. 10, 2022.

\bibitem{lagouvardos2020}
S. Lagouvardos, J. Dolby, N. Grech, A. Antoniadis, and Y. Smaragdakis, ``Static analysis of shape in
TensorFlow programs,'' in \emph{Proc. 34th European Conf. on Object-Oriented Programming (ECOOP)},
LIPIcs vol. 166, 2020, pp. 15:1--15:29.

\bibitem{grafberger2022}
S. Grafberger, P. Groth, J. Stoyanovich, and S. Schelter, ``Data distribution debugging in machine
learning pipelines,'' \emph{The VLDB Journal}, vol. 31, pp. 1103--1126, 2022.

\bibitem{gerstl2026}
O. Gerstl, F. Sihler, and M. Tichy, ``Inferring the shape of data frames in R programs using
abstract interpretation,'' arXiv:2607.03889, 2026.

\bibitem{sihler2026}
F. Sihler, L. Pfrenger, O. Gerstl, and M. Tichy, ``Towards automatically inferring constraints to
identify implicit assumptions in data analysis,'' arXiv:2607.03114, 2026.

\bibitem{mirji2026}
V. Mirji, ``Shift schema drift left: Policy-aware compile-time contracts for typed JVM and Spark
pipelines,'' arXiv:2604.16986, 2026.

\bibitem{petricek2016}
T. Petricek, G. Guerra, and D. Syme, ``Types from data: Making structured data first-class citizens
in F\#,'' in \emph{Proc. 37th ACM SIGPLAN Conf. on Programming Language Design and Implementation
(PLDI)}, 2016, pp. 477--490.

\bibitem{flanagan2006}
C. Flanagan, ``Hybrid type checking,'' in \emph{Proc. 33rd ACM SIGPLAN-SIGACT Symp. on Principles
of Programming Languages (POPL)}, 2006, pp. 245--256.

\bibitem{cartwright1991}
R. Cartwright and M. Fagan, ``Soft typing,'' in \emph{Proc. ACM SIGPLAN Conf. on Programming
Language Design and Implementation (PLDI)}, 1991, pp. 278--292.

\bibitem{cousot1977}
P. Cousot and R. Cousot, ``Abstract interpretation: A unified lattice model for static analysis of
programs by construction or approximation of fixpoints,'' in \emph{Proc. 4th ACM Symp. on
Principles of Programming Languages (POPL)}, 1977, pp. 238--252.

\bibitem{cousot1976}
P. Cousot and R. Cousot, ``Static determination of dynamic properties of programs,'' in \emph{Proc.
2nd Int. Symp. on Programming}, Paris, 1976, pp. 106--130.

\bibitem{codd1979}
E. F. Codd, ``Extending the database relational model to capture more meaning,'' \emph{ACM Trans.
Database Systems}, vol. 4, no. 4, pp. 397--434, 1979.

\bibitem{moerkotte1998}
G. Moerkotte, ``Small materialized aggregates: A light weight index structure for data
warehousing,'' in \emph{Proc. 24th Int. Conf. on Very Large Data Bases (VLDB)}, 1998, pp. 476--487.

\bibitem{melnik2010}
S. Melnik, A. Gubarev, J. J. Long, G. Romer, S. Shivakumar, M. Tolton, and T. Vassilakis, ``Dremel:
Interactive analysis of web-scale datasets,'' \emph{Proc. VLDB Endowment}, vol. 3, no. 1--2,
pp. 330--339, 2010.

\bibitem{parquet}
Apache Parquet, ``File format,'' \url{https://parquet.apache.org/docs/file-format/}.

\bibitem{ragan2013}
J. Ragan-Kelley, C. Barnes, A. Adams, S. Paris, F. Durand, and S. Amarasinghe, ``Halide: A language
and compiler for optimizing parallelism, locality, and recomputation in image processing
pipelines,'' in \emph{Proc. 34th ACM SIGPLAN Conf. on Programming Language Design and
Implementation (PLDI)}, 2013, pp. 519--530.

\bibitem{palkar2017}
S. Palkar, J. J. Thomas, A. Shanbhag, D. Narayanan, H. Pirk, M. Schwarzkopf, S. Amarasinghe, and
M. Zaharia, ``Weld: A common runtime for high performance data analytics,'' in \emph{Proc. Conf. on
Innovative Data Systems Research (CIDR)}, 2017.

\bibitem{jia2011}
Y. Jia and M. Harman, ``An analysis and survey of the development of mutation testing,'' \emph{IEEE
Trans. Software Engineering}, vol. 37, no. 5, pp. 649--678, 2011.

\bibitem{islam2019}
M. J. Islam, G. Nguyen, R. Pan, and H. Rajan, ``A comprehensive study on deep learning bug
characteristics,'' in \emph{Proc. 27th ACM Joint Meeting on European Software Engineering Conf. and
Symp. on the Foundations of Software Engineering (ESEC/FSE)}, 2019, pp. 510--520.

\bibitem{humbatova2020}
N. Humbatova, G. Jahangirova, G. Bavota, V. Riccio, A. Stocco, and P. Tonella, ``Taxonomy of real
faults in deep learning systems,'' in \emph{Proc. 42nd Int. Conf. on Software Engineering (ICSE)},
2020, pp. 1110--1121.

\bibitem{huynh2008}
Q. Huynh-Thu and M. Ghanbari, ``Scope of validity of PSNR in image/video quality assessment,''
\emph{Electronics Letters}, vol. 44, no. 13, pp. 800--801, 2008.

\bibitem{necula1997}
G. C. Necula, ``Proof-carrying code,'' in \emph{Proc. 24th ACM SIGPLAN-SIGACT Symp. on Principles of
Programming Languages (POPL)}, 1997, pp. 106--119.

\bibitem{holzle1992}
U. H\"olzle, C. Chambers, and D. Ungar, ``Debugging optimized code with dynamic deoptimization,'' in
\emph{Proc. ACM SIGPLAN Conf. on Programming Language Design and Implementation (PLDI)}, 1992,
pp. 32--43.

\bibitem{markl2004}
V. Markl, V. Raman, D. Simmen, G. Lohman, H. Pirahesh, and M. Cilimdzic, ``Robust query processing
through progressive optimization,'' in \emph{Proc. ACM SIGMOD Int. Conf. on Management of Data},
2004, pp. 659--670.

\bibitem{ioannidis1992}
Y. E. Ioannidis, R. T. Ng, K. Shim, and T. K. Sellis, ``Parametric query optimization,'' in
\emph{Proc. 18th Int. Conf. on Very Large Data Bases (VLDB)}, 1992, pp. 103--114.

\bibitem{cousot2013}
P. Cousot, R. Cousot, M. F\"ahndrich, and F. Logozzo, ``Automatic inference of necessary
preconditions,'' in \emph{Proc. 14th Int. Conf. on Verification, Model Checking, and Abstract
Interpretation (VMCAI)}, LNCS 7737, 2013, pp. 128--148.

\bibitem{ahmed2024}
S. Ahmed, H. Gao, and H. Rajan, ``Inferring data preconditions from deep learning models for
trustworthy prediction in deployment,'' in \emph{Proc. 46th Int. Conf. on Software Engineering
(ICSE)}, 2024.

\bibitem{clinger1990}
W. D. Clinger, ``How to read floating point numbers accurately,'' in \emph{Proc. ACM SIGPLAN Conf.
on Programming Language Design and Implementation (PLDI)}, 1990, pp. 92--101.

\end{thebibliography}

\end{document}
